\documentclass[10pt]{article}
% Shared preamble for model_card.tex and assumptions_register.tex
\usepackage[T1]{fontenc}
\usepackage[utf8]{inputenc}
\usepackage{lmodern}
\usepackage[margin=2cm]{geometry}
\usepackage{amsmath,amssymb}
\usepackage{graphicx}
\usepackage{booktabs,longtable,tabularx,array,multirow}
\usepackage[table]{xcolor}
\usepackage{enumitem}
\usepackage{caption}
\usepackage[hidelinks]{hyperref}
\usepackage{xurl}
\usepackage{pdflscape}

\graphicspath{{figures/}}
\setlist{itemsep=2pt,topsep=3pt}
\captionsetup{font=small,labelfont=bf}
\renewcommand{\arraystretch}{1.15}

% ---- editing helpers -------------------------------------------------------
% \TODO{...} marks text the students still have to write or confirm (renders red in the PDF).
\newcommand{\TODO}[1]{\textcolor{red!75!black}{\textbf{[TODO:} #1\textbf{]}}}
% Single source for facts that appear in both documents; edit here, not in the text.
\newcommand{\Nall}{138{,}554}
\newcommand{\Nhh}{85{,}182}
\newcommand{\prev}{0.631}
\newcommand{\capacity}{30\%}
\newcommand{\epsTol}{0.10}
\newcommand{\minSupport}{30}
\newcommand{\nbDate}{25 October 2026}
% Team facts used by every document in writeup/, report/ and slides/. Edit here once.
\newcommand{\StudentA}{Sumedh Navuda}
\newcommand{\StudentB}{Aldino}
\newcommand{\Authors}{\StudentA{} and \StudentB{}}
\newcommand{\Contact}{sumedhnavuda@outlook.com}
\newcommand{\RollA}{[roll no.]}   % <- replace with Sumedh's roll number
\newcommand{\RollB}{[roll no.]}   % <- replace with Aldino's roll number
\newcommand{\RepoURL}{https://github.com/sumedhnvda/ict5143-p8-fairness-audits}
\newcommand{\TeamID}{[team ID]}   % <- replace with the team ID given by the instructor
\newcommand{\SubmitTag}{submission-v1}
   % names, roll numbers, contact, repository (shared with report/ and slides/)
\newcommand{\nb}{\texttt{Problem8\_\allowbreak Intersectional\_\allowbreak Fairness.ipynb}}
\newcommand{\sect}[1]{notebook~\S#1}
\newcommand{\status}[1]{\textsf{\small #1}}

\newcolumntype{L}[1]{>{\raggedright\arraybackslash}p{#1}}

\title{Literature Review and Research Positioning\\[2pt]\large ICT 5143 Capstone, Problem 8: reliable intersectional fairness audits for sparse public-coverage groups}
\author{\Authors}
\date{Version 1.0 \textperiodcentered{} \nbDate}

\begin{document}
\maketitle

\begin{abstract}
We review 19 papers on intersectional and subgroup fairness, estimating fairness metrics for small groups, statistical
inference for audits, generalisation of fairness constraints, benchmark data and labels, and robust hierarchical shrinkage.
The review answers the handout's gap-verification questions. Hierarchical Bayesian estimation of intersectional fairness
already exists (Foulds et al.~\cite{foulds2020bayes}), as does attribute-structured estimation for small subgroups
(Herlihy et al.~\cite{herlihy2024}). Contamination and heavy-tailed priors that let outlying units escape shrinkage are
classical (Berger and Berliner~\cite{berger1986}; Datta and Lahiri~\cite{datta1995}). We therefore narrow the
contribution to an \emph{evaluation}: when Bayesian pooling produces false ``$\Delta\text{EOpp}\le\epsilon$'' conclusions
for sparse intersections, including validation-time rule selection. We state it as one falsifiable sentence (H1) and
test it out of sample.
\end{abstract}

\section{Search method}
We searched arXiv, the ACM Digital Library (FAccT), PMLR (ICML, AISTATS), the NeurIPS and IJCAI proceedings, SIAM (SDM),
Project Euclid and journal sites. Search terms combined \emph{intersectional fairness}, \emph{subgroup fairness},
\emph{small groups}, \emph{Bayesian}, \emph{hierarchical}, \emph{shrinkage}, \emph{fairness audit},
\emph{uncertainty}, \emph{contamination prior} and \emph{robust small area estimation}. We started from the references
in the handout (Foulds et al., Kearns et al., Ding et al., Folktables) and followed citations. A paper was included if it
(i) defines intersectional or subgroup fairness, (ii) estimates or tests fairness metrics under small group sizes,
(iii) studies whether fairness guarantees generalise, (iv) supplies our data or label caveats, or (v) supplies the
shrinkage method we modify. Every bibliographic entry was checked against its published abstract or proceedings page.
For several papers only the abstract was available to us. The table therefore records only what the abstract states and
marks other details ``not stated in abstract''. This is a limitation of the review: rows built from abstracts may omit assumptions stated only in the
full text. The closest papers (\S\ref{sec:closest}) are the ones whose full text matters most.

\section{Thematic review}
\subsection{Intersectional and subgroup fairness definitions}
Kearns et al.~\cite{kearns2018} show that fairness over a few pre-defined groups can hide unfairness to structured
subgroups (``fairness gerrymandering''). They prove that auditing subgroup fairness is computationally equivalent to weak
agnostic learning, and they give a learner--auditor game. H\'ebert-Johnson et al.~\cite{hebert2018} require calibration on
every computationally identifiable subpopulation (multicalibration). Foulds et al.~\cite{foulds2020icde} propose
\emph{differential fairness}, an intersectional definition motivated by differential privacy. Molina and
Loiseau~\cite{molina2022} characterise when marginal fairness bounds intersectional fairness, and report that the bounds
are loose in practice. Wang et al.~\cite{wang2022} discuss which identities to include, how to handle shrinking
subgroups, and how to evaluate more subgroups, with experiments on US Census tabular data. Gohar and Cheng~\cite{gohar2023}
survey intersectional notions and mitigation. \emph{Relevance:} our Experiment~1 (single attribute vs intersectional) and
our interaction-only semi-synthetic world are instances of the phenomenon these papers establish. They are not new.

\subsection{Estimating fairness metrics for small groups}
Foulds et al.~\cite{foulds2020bayes}, the reference the handout names as establishing the gap, show that intersectional
fairness measurement becomes unreliable as data become sparse. They propose Bayesian models with different bias--variance
trade-offs, including a new hierarchical Bayesian logistic regression, and Bayesian model averaging, evaluated on criminal
justice, census and synthetic data. Ji et al.~\cite{ji2020} estimate group fairness metrics with a hierarchical Bayesian
calibration model that also uses unlabeled data, and report lower estimation error. Morina et al.~\cite{morina2019} extend
single-attribute metrics to intersections and describe estimation when subgroups have little data. Lum et
al.~\cite{lum2022} show that naive between-group disparity estimates are statistically biased and propose a
``double-corrected'' variance estimator. After correction, apparent disparities were no longer significant. Herlihy et
al.~\cite{herlihy2024} use structured regression to obtain reliable performance estimates and confidence intervals for
very small intersectional subgroups, with goodness-of-fit tests that reveal the structure of harms. Himmelreich et
al.~\cite{himmelreich2025} argue that the intersectionality problem is one of statistical uncertainty due to small data,
and sketch a hypothesis-testing proposal. \emph{Relevance:} uncertainty-aware intersectional estimation, Bayesian or
structured, is established. The upward bias of raw max$-$min gaps that we see in simulation agrees with Lum et al.

\subsection{Statistical inference for fairness audits}
Cherian and Cand\`es~\cite{cherian2024} cast auditing as multiple hypothesis testing and use the bootstrap to bound
performance gaps over many subgroups at once, with guarantees, either flagging or certifying subgroups. \emph{Relevance:}
this is the frequentist counterpart of certifying ``$\Delta\le\epsilon$''. We do not reproduce it (see limitations).

\subsection{Fairness criteria and generalisation of fairness guarantees}
Hardt et al.~\cite{hardt2016} define equality of opportunity and give post-processing threshold adjustment. Our
$\lambda=1$ rule is the ordinary opportunity-constrained baseline. Cotter et al.~\cite{cotter2019} show that
data-dependent fairness constraints satisfied during training need not hold at evaluation, and they enforce constraints on
a separate validation set. \emph{Relevance:} the general failure of fairness constraints to transfer is known. Our
narrower question is whether \emph{pooled Bayesian audit criteria} make validation-time tolerance claims that fail on
independent test data, and whether they do so more often than raw or independent criteria.

\subsection{Benchmark data and label validity}
Ding et al.~\cite{ding2021} build the Folktables prediction tasks from US Census data across years and states. We use
the ACSPublicCoverage task through the Folktables package~\cite{folktables}. Fogliato et al.~\cite{fogliato2020} show
that small group-dependent label biases can overturn fairness conclusions. \emph{Relevance:} our label is
\emph{absence of public coverage}, not uninsurance; label noise is listed as an untested assumption (D2 in the register).

\subsection{Robust hierarchical shrinkage}
Berger and Berliner~\cite{berger1986} study $\epsilon$-contaminated prior classes and type-II maximum-likelihood
(empirical-Bayes) selection of a prior from the class. Datta and Lahiri~\cite{datta1995} note that normal-prior Bayes
estimators perform poorly with even one extreme small area, and use scale mixtures of normals (e.g.\ Cauchy for outlying
areas) for robust hierarchical Bayes small-area estimation. \emph{Relevance:} letting an atypical group escape shrinkage
through a contamination or heavy-tailed prior is classical. \textbf{Our contamination prior is an application of this
idea, not a new method.}

\section{Comparison table}
Table~\ref{tab:comparison} compares the papers on the six dimensions required by the handout.

\begin{landscape}
\footnotesize
\begin{longtable}{@{}L{2.9cm}L{3.2cm}L{2.5cm}L{3.3cm}L{2.8cm}L{2.9cm}L{3.6cm}@{}}
\caption{Comparison of related work. ``n.s.'' = not stated in the abstract we verified.}\label{tab:comparison}\\
\toprule
Paper & Key assumptions & Data & Graphical / statistical model & Fairness definition & Inference & Limitation for our question \\
\midrule\endfirsthead
\toprule
Paper & Key assumptions & Data & Graphical / statistical model & Fairness definition & Inference & Limitation for our question \\
\midrule\endhead
\bottomrule\endfoot
Foulds et al.\ 2020 (SDM)~\cite{foulds2020bayes} & Sparse intersections can borrow strength through a probabilistic model & Criminal justice, census, synthetic & Four Bayesian models incl.\ hierarchical Bayesian logistic regression; model averaging & Intersectional fairness (differential fairness~\cite{foulds2020icde}) & Bayesian posterior; n.s.\ beyond that & Outcome-rate fairness, not TPR-based EOpp; masking of an exceptional group and validation-time selection not studied (n.s.) \\
Foulds et al.\ 2020 (ICDE)~\cite{foulds2020icde} & Protected attributes known; a privacy-style ratio bound is a meaningful fairness notion & COMPAS, census & -- & Differential fairness ($\epsilon$-DF) & Learning algorithm satisfying DF & Definition paper; no small-group uncertainty study \\
Kearns et al.\ 2018~\cite{kearns2018} & Subgroups from a structured function class & n.s.\ (primarily theory) & -- & Subgroup fairness over rich classes & Learner--auditor zero-sum game; agnostic-learning oracle & Auditing is hard in the worst case; finite-sample small-group estimation not the focus \\
H\'ebert-Johnson et al.\ 2018~\cite{hebert2018} & Subpopulations computationally identifiable & n.s.\ (theory) & -- & Multicalibration & Algorithms with complexity analysis & Calibration criterion, not EOpp tolerance audits \\
Molina \& Loiseau 2022~\cite{molina2022} & Relations between marginal and intersectional fairness & n.s. & -- & Intersectional vs marginal fairness & High-probability bounds; grouping heuristic & Bounds loose in practice (authors) \\
Wang et al.\ 2022~\cite{wang2022} & Including more identities creates small subgroups & US Census tabular & -- & New intersectional evaluation metrics & Empirical & Pipeline guidance; no Bayesian audit calibration \\
Gohar \& Cheng 2023~\cite{gohar2023} & -- (survey) & -- & -- & Taxonomy of intersectional notions & -- & Survey \\
Ji et al.\ 2020~\cite{ji2020} & Many unlabeled, few labeled examples; scores calibratable per group & Several fairness datasets & Hierarchical latent-variable calibration model & Several group metrics & Bayesian posterior & Groups, not sparse intersections; no tolerance-decision error rates \\
Morina et al.\ 2019~\cite{morina2019} & Single-attribute metrics extend to intersections & One real dataset & n.s. & Intersectional extensions of standard metrics & Estimation for small subgroups (n.s.); post-processing & No pooling-masking analysis \\
Lum et al.\ 2022~\cite{lum2022} & Group metrics carry sampling noise & Simulation and one real dataset & Variance decomposition (frequentist) & Variance of performance across groups & Double-corrected variance estimator & Summary of spread; not max$-$min tolerance certification \\
Herlihy et al.\ 2024~\cite{herlihy2024} & Subgroup performance has regression structure in group features & Two public datasets, semi-synthetic variants & Structured regression & Subgroup performance metrics & Confidence intervals; goodness-of-fit tests & Closest structured method; false-tolerance rates under an isolated interaction n.s. \\
Himmelreich et al.\ 2025~\cite{himmelreich2025} & Intersectional problem = statistical uncertainty & Partly empirical & -- & Desiderata for intersectional fairness & Simple hypothesis testing (sketch) & Conceptual; no calibrated procedure \\
Cherian \& Cand\`es 2024~\cite{cherian2024} & Model-agnostic; bootstrap validity & Benchmark datasets & -- & Any metric / group criterion & Simultaneous bootstrap bounds (multiple testing) & Frequentist; does not compare Bayesian pooling \\
Hardt et al.\ 2016~\cite{hardt2016} & Group-conditional rates estimable & n.s. & -- & Equal opportunity; equalised odds & Post-processing threshold adjustment & Small-group estimation error not addressed \\
Cotter et al.\ 2019~\cite{cotter2019} & Constraints can be enforced on held-out data & n.s. & -- & Data-dependent fairness constraints & Two-player constrained optimisation & Training-time generalisation, not audit claims \\
Ding et al.\ 2021~\cite{ding2021} & Census-derived tasks represent the population & ACS PUMS, all states, multiple years & -- & (benchmark) & -- & Dataset; ours uses its ACSPublicCoverage task \\
Fogliato et al.\ 2020~\cite{fogliato2020} & Observed label is a noisy proxy & Two criminal justice datasets & Sensitivity-analysis framework & Predictive bias & Sensitivity analysis & Different domain; motivates label caveat \\
Berger \& Berliner 1986~\cite{berger1986} & Base prior plus $\epsilon$-contamination class & -- (theory) & $\epsilon$-contaminated priors & -- & Posterior ranges; type-II ML & Not about fairness \\
Datta \& Lahiri 1995~\cite{datta1995} & Some small areas are outliers & n.s. & Hierarchical Bayes with scale-mixture-of-normal priors & -- & Robust hierarchical Bayes & Not about fairness or tolerance claims \\
\midrule
\textbf{This project} & Exchangeable or attribute-structured pooling of TPRs; binomial counts & ACS PUMS 2018 CA (core), TX (replication) & Discrete BN (predictor); beta-binomial, hierarchical logistic, contamination audit models & Intersectional EOpp (primary), DP (secondary) & VE / clique-tree; EB, grid full Bayes, quadrature; simulation calibration & EB plug-in under-covers; calibration design-specific; struct not an exact reproduction of~\cite{foulds2020bayes} \\
\end{longtable}
\end{landscape}

\section{Closest existing method and baseline reproduction}\label{sec:closest}
The closest methods are the hierarchical Bayesian logistic regression of Foulds et al.~\cite{foulds2020bayes} and the
structured regression of Herlihy et al.~\cite{herlihy2024}. Both let a small intersection borrow strength through its
attributes. As the handout allows, we reproduce a \emph{clearly documented restricted version} (\texttt{struct}, notebook
\S6):
\[
\operatorname{logit}\theta_g=\beta_0+\beta_{\text{race}[g]}+\beta_{\text{sex}[g]}+\beta_{\text{dis}[g]}+u_g,\qquad
u_g\sim N(0,\tau^2),\qquad k_g\sim\text{Bin}(n_g,\theta_g),
\]
applied to true-positive rates. $\beta$ and $\tau$ are fitted by marginal likelihood (adaptive quadrature), and $u_g$ is
drawn from its grid posterior. Differences from Foulds et al.\ as far as we could verify: we model TPRs among positives
rather than outcome rates; we use only main effects; we plug in empirical-Bayes hyperparameters; and we did not match their
priors, covariates or inference algorithm. Exact reproduction is therefore \textbf{not} claimed.
An exact reproduction (their code, priors and sampler on our counts) was not attempted and remains future work.
The quadrature is checked against brute-force integration ($10^{-10}$). Under its own generative model, its 90\% intervals
cover 0.665, against 0.905 for the same sampler with the true hyperparameters (notebook \S6.1).

\section{Gap verification and narrowing of the claim}
\begin{tabularx}{\textwidth}{@{}L{4.6cm}X@{}}
\toprule
Candidate gap (handout) & What the literature already answers, and what remains \\
\midrule
Hierarchical Bayesian intersectional estimation & Established~\cite{foulds2020bayes,ji2020,herlihy2024}. Not claimed. \\
Small groups make fairness metrics unreliable & Established~\cite{foulds2020bayes,lum2022,himmelreich2025}. Not claimed. \\
Marginal audits can miss intersectional harm & Established~\cite{kearns2018,molina2022}. Reported as a replication on our data (Exp.~1). \\
A prior that lets atypical groups escape shrinkage & Classical in robust and small-area Bayes~\cite{berger1986,datta1995}. Claimed only as an \emph{application}, evaluated as an audit safeguard. We found no application to fairness audits; this search is not exhaustive. \\
Fairness guarantees may not generalise from training to test & Established for training constraints~\cite{cotter2019}. \textbf{Remaining:} whether \emph{pooled audit criteria} produce validation-time tolerance claims that fail on test (H1), and how this compares with raw and independent criteria. \\
Pooling masks a small exceptional group; quantile rules as a fix; calibrating a tolerance procedure & Not answered in the papers we reviewed. We add a diagnosis and a tested fix (design-matched calibration, held-out certification). \textbf{Remaining:} a controlled evaluation of false-tolerance rates by deficit structure (cell, interaction-only, main effect), thinning and design, and whether a simulation calibration transfers to the deployment design. \\
\bottomrule
\end{tabularx}

\medskip
\noindent\textbf{Narrowed contribution.} A controlled and real-data evaluation of when Bayesian pooling, exchangeable and
attribute-structured, produces false ``$\Delta\text{EOpp}\le\epsilon$'' conclusions for sparse intersections. It covers
minority thinning, interaction-only and main-effect deficits, validation-time rule selection, and the transfer of a
simulation-calibrated robust-prior procedure between designs.

\section{Falsifiable contribution}
\textbf{H1.} \emph{When the group-threshold rule's $\lambda$ is selected on validation data and a pooled Bayesian audit
criterion (hierarchical, full-Bayes, complete, attribute-structured or contamination pooling) reports intersectional
$\Delta\text{EOpp}\le0.10$, the frozen rule's empirical test $\Delta\text{EOpp}$ exceeds 0.10 in more than half of those cases.}

H1 was formed from the California results and tested out of sample on Texas 2018 (notebook \S14.1), with the decision
rule fixed before the test was run.
% generated by make_tables.py from ../results; do not edit by hand
\begin{center}\begin{tabular}{lc}
\toprule
Texas 2018, 5 household splits & value \\
\midrule
pooled selections & 40 \\
validation-time claims $\Delta\text{EOpp}\le0.10$ (pooled) & 27 \\
claims failing on test ($\Delta\text{EOpp}>0.10$) & 27 \\
test $\Delta$EOpp of the claims & 0.110--0.888 \\
claims by raw / independent criteria & 0 \\
\midrule
\textbf{H1 verdict} & \textbf{supported} \\
\bottomrule
\end{tabular}\end{center}


\noindent\textbf{Methodological contribution and its test (H2).} The review found no application of robust priors to
fairness audits and no study of calibration transfer between audit designs. We therefore propose \emph{design-matched
calibration} plus \emph{held-out certification} (notebook \S15) and test it with preregistered hypotheses on the Texas
design. Its audit score $J=\text{ACC}-\text{FTS}$ is 0.664, against at most 0.385 for every established audit (H2a
supported) and 0.189 for the invented-design calibration (H2b supported). Held-out certification produced no false
tolerance claims on real CA or TX data, against 10 of 10 with same-data certification (H2c supported).

\noindent\textbf{Earlier draft claim.} The draft claim ``an $\epsilon$-contamination prior reduces the false-satisfaction
rate of an empirical-Bayes upper-quantile audit by at least half'' is not supported for the calibrated procedure
(0.980 $\to$ 0.503; the required value is $\le0.490$).

% Shared bibliography (\input by literature_review.tex and manuscript.tex). Every entry verified against its
% published abstract or proceedings page.
\begin{thebibliography}{99}\small
\bibitem{foulds2020bayes} J.~R. Foulds, R.~Islam, K.~N. Keya, S.~Pan. Bayesian modeling of intersectional fairness: the variance of bias. \emph{Proc.\ SIAM Int.\ Conf.\ Data Mining (SDM)}, 424--432, 2020. doi:10.1137/1.9781611976236.48. arXiv:1811.07255.
\bibitem{foulds2020icde} J.~R. Foulds, R.~Islam, K.~N. Keya, S.~Pan. An intersectional definition of fairness. \emph{Proc.\ IEEE 36th Int.\ Conf.\ Data Engineering (ICDE)}, 1918--1921, 2020. doi:10.1109/ICDE48307.2020.00203.
\bibitem{kearns2018} M.~Kearns, S.~Neel, A.~Roth, Z.~S. Wu. Preventing fairness gerrymandering: auditing and learning for subgroup fairness. \emph{Proc.\ ICML}, PMLR 80:2564--2572, 2018.
\bibitem{hebert2018} \'U.~H\'ebert-Johnson, M.~P. Kim, O.~Reingold, G.~N. Rothblum. Multicalibration: calibration for the (computationally-identifiable) masses. \emph{Proc.\ ICML}, PMLR 80:1939--1948, 2018.
\bibitem{molina2022} M.~Molina, P.~Loiseau. Bounding and approximating intersectional fairness through marginal fairness. \emph{Advances in NeurIPS 35}, 2022. arXiv:2206.05828.
\bibitem{wang2022} A.~Wang, V.~V. Ramaswamy, O.~Russakovsky. Towards intersectionality in machine learning: including more identities, handling underrepresentation, and performing evaluation. \emph{Proc.\ ACM FAccT}, 336--349, 2022. doi:10.1145/3531146.3533101.
\bibitem{gohar2023} U.~Gohar, L.~Cheng. A survey on intersectional fairness in machine learning: notions, mitigation, and challenges. \emph{Proc.\ IJCAI} (Survey Track), 6619--6627, 2023. doi:10.24963/ijcai.2023/742.
\bibitem{ji2020} D.~Ji, P.~Smyth, M.~Steyvers. Can I trust my fairness metric? Assessing fairness with unlabeled data and Bayesian inference. \emph{Advances in NeurIPS 33}, 18600--18612, 2020. arXiv:2010.09851.
\bibitem{morina2019} G.~Morina, V.~Oliinyk, J.~Waton, I.~Maru\v{s}i\'c, K.~Georgatzis. Auditing and achieving intersectional fairness in classification problems. arXiv:1911.01468, 2019.
\bibitem{lum2022} K.~Lum, Y.~Zhang, A.~Bower. De-biasing ``bias'' measurement. \emph{Proc.\ ACM FAccT}, 379--389, 2022. doi:10.1145/3531146.3533105.
\bibitem{herlihy2024} C.~Herlihy, K.~Truong, A.~Chouldechova, M.~Dud\'ik. A structured regression approach for evaluating model performance across intersectional subgroups. \emph{Proc.\ ACM FAccT}, 2024. arXiv:2401.14893.
\bibitem{himmelreich2025} J.~Himmelreich, A.~Hsu, E.~Veomett, K.~Lum. The intersectionality problem for algorithmic fairness. \emph{Algorithmic Fairness Through the Lens of Metrics and Evaluation}, PMLR 279:68--95, 2025. arXiv:2411.02569.
\bibitem{cherian2024} J.~J. Cherian, E.~J. Cand\`es. Statistical inference for fairness auditing. \emph{Journal of Machine Learning Research} 25(149):1--49, 2024.
\bibitem{hardt2016} M.~Hardt, E.~Price, N.~Srebro. Equality of opportunity in supervised learning. \emph{Advances in NIPS 29}, 3315--3323, 2016.
\bibitem{cotter2019} A.~Cotter, M.~Gupta, H.~Jiang, N.~Srebro, K.~Sridharan, S.~Wang, B.~Woodworth, S.~You. Training well-generalizing classifiers for fairness metrics and other data-dependent constraints. \emph{Proc.\ ICML}, PMLR 97:1397--1405, 2019.
\bibitem{ding2021} F.~Ding, M.~Hardt, J.~Miller, L.~Schmidt. Retiring Adult: new datasets for fair machine learning. \emph{Advances in NeurIPS 34}, 6478--6490, 2021. arXiv:2108.04884.
\bibitem{folktables} Folktables project. Dataset definitions, documentation and data-access tools. \url{https://github.com/socialfoundations/folktables}.
\bibitem{fogliato2020} R.~Fogliato, A.~Chouldechova, M.~G'Sell. Fairness evaluation in presence of biased noisy labels. \emph{Proc.\ AISTATS}, PMLR 108:2325--2336, 2020.
\bibitem{berger1986} J.~O. Berger, L.~M. Berliner. Robust Bayes and empirical Bayes analysis with $\epsilon$-contaminated priors. \emph{Annals of Statistics} 14(2):461--486, 1986. doi:10.1214/aos/1176349933.
\bibitem{datta1995} G.~S. Datta, P.~Lahiri. Robust hierarchical Bayes estimation of small area characteristics in the presence of covariates and outliers. \emph{Journal of Multivariate Analysis} 54(2):310--328, 1995.
\end{thebibliography}

\end{document}
